\begin{abstract}
Este trabalho aborda dois problemas de autovalores para matrizes reais construídas a partir da matrícula do aluno. No primeiro, calcula-se a decomposição real de Schur por meio da redução ortogonal à forma de Hessenberg, seguida de iterações QR com rotações de Givens. A partir da matriz de Schur, obtêm-se os autovalores e examina-se a sensibilidade por números de condição associados aos autovalores simples. No segundo problema, aplica-se o método da potência, com normalização na norma euclidiana e estimativa do autovalor pelo quociente de Rayleigh, para determinar o autovalor dominante. Para a matrícula 318375, o primeiro algoritmo convergiu em 68 iterações, produzindo autovalores reais \(26{,}68556904\), \(18{,}42664022\), \(2{,}88779074\) e \(34\). O método da potência convergiu em 818 iterações para \(10^{-8}\), fornecendo \(\lambda_{\mathrm{dom}}\approx17{,}6085936704\). Verificações de ortogonalidade, reconstrução, invariância do traço e resíduo do par próprio foram utilizadas para avaliar os resultados. Observou-se ainda que a primeira matriz é não normal e que seus autovalores apresentam sensibilidades distintas.
\end{abstract}

\noindent\textbf{Palavras-chave:} problema de autovalores; decomposição de Schur; redução de Hessenberg; rotações de Givens; método da potência; sensibilidade.
